% Generated by code/make_standalone.py. Edit the modular files.
\documentclass[11pt,a4paper]{article}
% BEGIN preamble.tex
\usepackage[T1]{fontenc}
\usepackage{lmodern,microtype}
\usepackage[a4paper,margin=26mm,headheight=15pt]{geometry}
\usepackage{amsmath,amssymb,amsthm,mathtools}
\usepackage{mathrsfs}
\usepackage{booktabs,longtable,array,enumitem}
\usepackage[dvipsnames]{xcolor}
\usepackage{fancyhdr}
\usepackage[colorlinks=true,linkcolor=MidnightBlue,citecolor=MidnightBlue,urlcolor=MidnightBlue,breaklinks=true]{hyperref}
\usepackage{bookmark}
\pagestyle{fancy}\fancyhf{}
\lhead{\small Commensurate Lerch calculus}\rhead{\small ProveIt research continuation}
\cfoot{\thepage}\renewcommand{\headrulewidth}{0.3pt}
\numberwithin{equation}{section}
\newtheorem{theorem}{Theorem}[section]
\newtheorem{lemma}[theorem]{Lemma}
\newtheorem{proposition}[theorem]{Proposition}
\newtheorem{corollary}[theorem]{Corollary}
\theoremstyle{definition}\newtheorem{definition}[theorem]{Definition}
\theoremstyle{remark}\newtheorem{remark}[theorem]{Remark}
\DeclareMathOperator{\Li}{Li}
\DeclareMathOperator{\Rea}{Re}
\DeclareMathOperator{\Ima}{Im}
\DeclareMathOperator{\Res}{Res}
\newcommand{\dd}{\,\mathrm d}
\newcommand{\C}{\mathbb C}\newcommand{\R}{\mathbb R}\newcommand{\N}{\mathbb N}
\newcommand{\bq}{\boldsymbol q}\newcommand{\bs}{\boldsymbol s}
\newcommand{\bA}{\boldsymbol A}\newcommand{\bdelta}{\boldsymbol\delta}
\newcommand{\CC}{\mathcal C}\newcommand{\UU}{\mathcal U}\newcommand{\EE}{\mathcal E}
\newcommand{\HH}{\mathcal H}\newcommand{\PP}{\mathcal P}
\newcommand{\src}[1]{\nolinkurl{#1}}
\setlist{itemsep=4pt,topsep=6pt}
\allowdisplaybreaks[2]
\hypersetup{pdftitle={Commensurate Lerch Calculus},pdfauthor={Research continuation for Vladimir Reshetnikov's ProveIt programme},pdfsubject={Finite cyclotomic closure, unequal-center order jets, harmonic remainders and Stieltjes--Gamma identities}}

% END preamble.tex
\title{\textbf{Commensurate Lerch Calculus}\\[6pt]
\large Finite cyclotomic closure, unequal-center expansions,\\
harmonic remainders, and Stieltjes--Gamma identities}
\author{Research continuation for Vladimir Reshetnikov's ProveIt programme\\
\small Mathematical derivations and computational checks with OpenAI assistance}
\date{10 October 2026}
\begin{document}
\maketitle
\begin{abstract}
We extend reciprocal Lerch convolution from equal logarithmic scales to
arbitrary commensurate positive scales, at every depth and every independent
complex order. After a necessary scale normalization and alignment of the
affine resolvents, the convolution depends only on the total order. A finite
pole-lattice calculation reduces it to Lerch functions, and at the central
point to pole-free Hurwitz-zeta combinations. We prove all mixed logarithmic
moment identities, derive all-index Stieltjes and log-Gamma consequences,
and give a normally convergent expansion for unequal affine centers,
including every independent order derivative.

The resulting two-scale family includes
$\int_0^\infty\log(1+y)\log(1+y^{-2})\,\dd y/y
=\pi G-3\zeta(3)/8$, balanced unequal harmonic-remainder products, and
explicit Gamma antiderivatives. A separate residue parity theorem works even
for incommensurate scales; for example every central three-factor logistic
convolution is $\pi^2\sum_jq_j^{-2}/12$. Two precise rigidity statements
identify when pure order addition and a common-step elementary kernel
representation are possible. These statements do not assert arithmetic
independence or unrestricted special-function nonreducibility.

The article supplies full analytic proofs, executable exact checks, and
numerical diagnostics. It resolves the commensurate-dilation and
unequal-shift-expansion research questions in the preceding reciprocal
report, in the explicitly stated domains. It does not resolve the canonical
Gaussian $S_6$ or revised $S_8$ conjectures, and makes no blanket claim of
historical priority for individual identities.
\end{abstract}
\vfill
\noindent\small\textbf{Source anchor.} Targeted canonical reading at repository commit\\
\src{338a3a3877c337f981072061c16bc8ee1d87155d}, together with the complete
accessible source of the preceding reciprocal report. The incoming ZIP
listing was inspected; its binary members were not audited. No repository
files were changed.\normalsize
\clearpage
\begingroup\small\tableofcontents\endgroup
\clearpage
% BEGIN sections/01-scope.tex
\section{Research target and the new exact families}
\label{clc:sec:scope}

The ProveIt manuscript already contains a substantial theory of
polylogarithmic identities, generalized harmonic sums, Hurwitz order
jets, Stieltjes distributions, and Gamma correlations \cite{Repo}.
The preceding \emph{Reciprocal Lerch Convolutions} report establishes
entire-order addition, finite central-factorial formulas, and all
logarithmic moments for equal logarithmic scales \cite{RLC}. Its proposed
research questions explicitly ask for commensurate dilations and for a
normally convergent unequal-shift expansion with independent order jets.
These are our principal targets. Re-proving an equal-scale identity is not
counted as a new result.

The central distinction is between a spectral order and a logarithmic
scale. If the argument $y$ is replaced by $y^q$, the bilateral transform
changes its cosecant period and also introduces a factor $q^s$. That
factor cannot be ignored. Once it is removed, the new periods can be
combined into a finite pole lattice exactly when their ratios are rational.
A separate expansion handles unequal affine centers without pretending
that a finite integer-order partial fraction identity can be differentiated
in its integer index.

\subsection{Two accessible consequences}
For arbitrary complex $s,t$, write $W=s+t$. One result proved below is
\begin{equation}\label{clc:eq:intro-beta}
 \boxed{\displaystyle
 \int_0^\infty \Li_s(-y)\Li_t(-y^{-2})\frac{\dd y}{y}
 =2^{t-1}\pi\beta(W)-2^{-s-1}W\eta(W+1).}
\end{equation}
Here $\beta$ is the Dirichlet beta function and $\eta$ the Dirichlet eta
function. Both sides are entire in the two orders, and the integral is
ordinary and absolutely convergent. In particular, with
$G=\beta(2)$ denoting Catalan's constant,
\begin{equation}\label{clc:eq:intro-log}
 \int_0^\infty \log(1+y)\log(1+y^{-2})\frac{\dd y}{y}
 =\pi G-\frac38\zeta(3).
\end{equation}
The two orders in \eqref{clc:eq:intro-beta} do not enter symmetrically.
The normalization restoring pure order addition is multiplication of the
integral by $2^{-t}$.

A different consequence does not require rationally related scales.
For any positive real $q_1,q_2,q_3$,
\begin{equation}\label{clc:eq:intro-cubic}
 \boxed{\displaystyle
 \int_{\R^2}
 \frac{1}{1+e^{-q_1x}}\,
 \frac{1}{1+e^{-q_2y}}\,
 \frac{1}{1+e^{q_3(x+y)}}\dd x\dd y
 =\frac{\pi^2}{12}\left(q_1^{-2}+q_2^{-2}+q_3^{-2}\right).}
\end{equation}
This is one case of an all-depth negative-integer parity theorem. It
illustrates that the absence of a common pole lattice does not rule out
exceptional exact central evaluations.

\subsection{Conventions}
We use the standard Hurwitz and Lerch functions \cite{DLMF25}:
\[
 \zeta(v,a)=\sum_{n=0}^\infty(n+a)^{-v},\qquad
 \Phi(z,v,a)=\sum_{n=0}^\infty\frac{z^n}{(n+a)^v},
\]
initially in their convergence domains and then by analytic continuation.
Throughout, $\Rea a>0$ unless another domain is explicitly stated; the
powers use the logarithm on the right half-plane. The rising factorial is
$(v)_j=v(v+1)\cdots(v+j-1)$, with $(v)_0=1$. A superscript $[h]$ means
an order derivative. A superscript $(h)$ on $\psi$ means an ordinary
argument derivative of the digamma function.

Our generalized Stieltjes convention is
\begin{equation}\label{clc:eq:stieltjes-convention}
 \zeta(1+v,a)=\frac1v+
 \sum_{n=0}^\infty\frac{(-1)^n}{n!}\gamma_n(a)v^n,
 \qquad \gamma_0(a)=-\psi(a).
\end{equation}
We repeatedly use the entire functions
\begin{align}
 \EE_a(v)&=v\zeta(v+1,a),&\EE_a(0)&=1,\label{clc:eq:E}\\
 \eta(v,a)&=\Phi(-1,v,a)
 =2^{-v}\left[\zeta\left(v,\frac a2\right)
                  -\zeta\left(v,\frac{a+1}2\right)\right].
 \label{clc:eq:eta}
\end{align}
In particular $\eta(v)=\eta(v,1)$, while
$\beta(v)=4^{-v}[\zeta(v,1/4)-\zeta(v,3/4)]$.
The singularities in these formulas are filled in before differentiation.

\subsection{Relationship to classical methods and remaining conjectures}
Euler's beta integral, Gamma reflection, convolution, pole summation,
and Bell-polynomial differentiation are classical tools; the use of real
convolution for Fermi--Dirac-type integrals is explicitly represented in
Selvaggi and Selvaggi \cite{SS}. We do not claim to have invented these
methods. The contribution here is the all-complex-order commensurate-scale
closure, its precise hypotheses, pole-free normalization, and the derived
families of identities and independent-center continuations.

The inspected canonical README records $S_4$ as proved and $S_6$ and the
revised $S_8$ candidate as conjectural. Those statuses are unchanged here.
A small residual is not an equality proof. The two research questions we
answer are the analytic continuation problems just described, not those
specific arithmetic reduction conjectures. Complete literature priority
for every special value has not been established.

% END sections/01-scope.tex
% BEGIN sections/02-analysis.tex
\section{Aligned profiles and entire-order addition}
\label{clc:sec:analysis}

For $y>0$ define
\begin{equation}\label{clc:eq:F}
 F_{a,s}(y)=y\Phi(-y,s,a),\qquad
 \ell(x)=\frac{e^x}{1+e^x}.
\end{equation}
Thus $F_{1,s}(y)=-\Li_s(-y)$ and $F_{a,0}(e^x)=\ell(x)$.
For a positive real scale $q$ and $\Rea A>-q$, set
\begin{equation}\label{clc:eq:h}
 h_{q,A,s}(x)=q^{-s}F_{1+A/q,s}(e^{qx}),\qquad
 g_q(p)=\frac{\pi}{q}\csc\frac{\pi p}{q}.
\end{equation}
We call $A$ the affine center. It is $q(a-1)$ in the original Lerch
shift, not $a$ itself.

\begin{lemma}[Weighted entire dependence]\label{clc:lem:weighted}
Let $q>0$, $\Rea A>-q$, and
\begin{equation}\label{clc:eq:strip}
 \max(0,-\Rea A)<c<q.
\end{equation}
Then $s\mapsto e^{-cx}h_{q,A,s}(x)$ is entire with values in
$L^1(\R)\cap L^\infty(\R)$. The assertion holds after any fixed
multiplication by a power of $x$, any fixed order or center derivative,
and locally uniformly in $A$. The profiles are smooth in $x$ and obey
\begin{equation}\label{clc:eq:profile-ladder}
 (\partial_x+A)h_{q,A,s}=h_{q,A,s-1}.
\end{equation}
\end{lemma}
\begin{proof}
For $\Rea s>0$, the Lerch integral and the change of variable from its
Laplace parameter to $qu$ give
\begin{equation}\label{clc:eq:fractional-profile}
 h_{q,A,s}(x)=\frac1{\Gamma(s)}
 \int_0^\infty u^{s-1}e^{-Au}\ell(q(x-u))\dd u.
\end{equation}
Every derivative of $\ell$ is $\ell$ times a polynomial in $\ell$.
Consequently $|\ell^{(j)}(v)|\le C_j\ell(v)$ on the real line.
The function $e^{-cx}\ell(qx)$ is bounded and integrable precisely for
$0<c<q$. Multiplying \eqref{clc:eq:fractional-profile} by $e^{-cx}$ and
writing $v=x-u$ separates the norm estimate into this logistic norm and
\[
 \int_0^\infty u^{\Rea s-1}e^{-(c+\Rea A)u}\dd u.
\]
Both are finite. On compact subsets of the initial parameter domain,
derivatives insert only powers of $u$, $\log u$, and smooth Gamma factors;
these are controlled by slightly larger integrable majorants.

Integration by parts in $u$, initially for $\Rea s>1$, proves
\eqref{clc:eq:profile-ladder}. For an arbitrary compact set of orders,
choose an integer $N$ with $\Rea(s+N)>0$ and define
$h_{q,A,s}=(\partial_x+A)^Nh_{q,A,s+N}$. The integration-by-parts identity
makes these definitions consistent on overlaps. They agree with the
convergent Lerch series for $x<0$ and hence with the specified continuation.
The derivative estimates for $\ell$ prove the norm-holomorphic extension.
Finally, choose neighboring admissible weights $c\pm\varepsilon$.
The bound $|x|^k\le C_{k,\varepsilon}e^{\varepsilon|x|}$ proves every
polynomially weighted assertion. Cauchy's formula in the parameters
supplies all fixed parameter derivatives.
\end{proof}

\begin{theorem}[Scaled bilateral transform]\label{clc:thm:transform}
Under \eqref{clc:eq:strip}, for every $s\in\C$,
\begin{equation}\label{clc:eq:transform}
 \int_\R e^{-px}h_{q,A,s}(x)\dd x
 =g_q(p)(p+A)^{-s},\qquad \Rea p=c.
\end{equation}
The power uses the right-half-plane logarithm of $p+A$.
\end{theorem}
\begin{proof}
For $\Rea s>0$, Fubini is justified by the preceding proof. The transform
of \eqref{clc:eq:fractional-profile} factors into $(p+A)^{-s}$ and
\[
 \int_\R e^{-px}\ell(qx)\dd x
 =\frac1q\int_0^\infty\frac{y^{-p/q}}{1+y}\dd y
 =\frac1q\Gamma(p/q)\Gamma(1-p/q)=g_q(p).
\]
The last step is Gamma reflection \cite{DLMF5}. Both sides are entire in
$s$, so the identity extends to all orders. In particular there is no
remaining $\Gamma(s)$ factor and no remaining $q^s$ factor.
\end{proof}

Let $\bq=(q_1,\ldots,q_r)$, $q_j>0$, and $q_*=\min_jq_j$.
For $r\ge2$ put $x_r=\mu-\sum_{j<r}x_j$ and define
\begin{equation}\label{clc:eq:C-def}
 \CC_{\bq,A}(W;\mu)=
 \int_{\R^{r-1}}\prod_{j=1}^r h_{q_j,A,s_j}(x_j)
           \dd x_1\cdots\dd x_{r-1},\qquad W=\sum_js_j.
\end{equation}
For $r=1$ the notation means evaluation of the single profile at $\mu$.
That the integral depends only on $W$ is part of the following theorem.

\begin{theorem}[Entire-order addition at arbitrary positive scales]
\label{clc:thm:addition}
If $\Rea A>-q_*$, \eqref{clc:eq:C-def} is absolutely convergent for all
independent complex orders, depends only on their sum, and is entire in
that sum. With
$G_{\bq}(p)=\prod_jg_{q_j}(p)$, it equals
\begin{equation}\label{clc:eq:C-contour}
 \CC_{\bq,A}(W;\mu)=\frac1{2\pi i}
 \int_{c-i\infty}^{c+i\infty}
 e^{p\mu}G_{\bq}(p)(p+A)^{-W}\dd p,
 \qquad \max(0,-\Rea A)<c<q_*.
\end{equation}
Every fixed mixed order derivative can be moved inside the ordinary
convolution integral. The contour expression is holomorphic for
$|\Ima\mu|<\pi\sum_jq_j^{-1}$.
\end{theorem}
\begin{proof}
Set $u_j(x)=e^{-cx}h_{q_j,A,s_j}(x)$. On the constraint hyperplane the
weight removed from the product is the constant $e^{c\mu}$. Thus the
absolute integral is at most
$e^{c\mu}\|u_r\|_\infty\prod_{j<r}\|u_j\|_1$.
Lemma~\ref{clc:lem:weighted} gives convergence and locally uniform bounds
for all derivatives. The product of Fourier transforms is
$G_{\bq}(p)(p+A)^{-W}$, by Theorem~\ref{clc:thm:transform}.
Along the vertical line this decays like
$\exp(-\pi|\Ima p|\sum_jq_j^{-1})$ times a fixed power of $1+|p|$
on compact parameter sets. Fourier inversion therefore applies. The
convolution is continuous, so equality holds pointwise. The same decay
proves the claimed complex $\mu$ strip and all holomorphic dependence.
\end{proof}

Write $K_{\bq}(\mu)=\CC_{\bq,A}(0;\mu)$; it is independent of $A$ and
is the convolution of the scaled logistic functions. An equivalent
representation, initially for $\Rea W>0$, is
\begin{equation}\label{clc:eq:base-fractional}
 \CC_{\bq,A}(W;\mu)=\frac1{\Gamma(W)}
 \int_0^\infty u^{W-1}e^{-Au}K_{\bq}(\mu-u)\dd u.
\end{equation}
It follows either by the Dirichlet simplex integral in
\eqref{clc:eq:fractional-profile} or by its bilateral transform. The
entire continuation in $W$, not the improper integral on the right,
is used when $\Rea W\le0$.

Scaling is also exact. For $\lambda>0$,
\begin{equation}\label{clc:eq:scaling}
 \CC_{\lambda\bq,\lambda A}(W;\mu)
 =\lambda^{1-r-W}\CC_{\bq,A}(W;\lambda\mu).
\end{equation}
Indeed $h_{\lambda q,\lambda A,s}(x)=\lambda^{-s}h_{q,A,s}(\lambda x)$,
and the $r-1$ integration variables supply the remaining power.

% END sections/02-analysis.tex
% BEGIN sections/03-pole-lattice.tex
\section{A finite pole lattice and pole-free Lerch closure}
\label{clc:sec:lattice}

Assume now that the positive scales are commensurate: there exists $Q>0$
with $d_j=Q/q_j\in\N$. For rational scales one may choose
\begin{equation}\label{clc:eq:Q}
 Q=\frac{\operatorname{lcm}(\operatorname{num}q_1,\ldots,
                              \operatorname{num}q_r)}
         {\gcd(\operatorname{den}q_1,\ldots,
                              \operatorname{den}q_r)}.
\end{equation}
This is a common multiple of the scales, not a common denominator.
Put
\begin{equation}\label{clc:eq:period-data}
 \sigma=(-1)^{\sum_jd_j},\qquad
 \PP=\bigcup_{j=1}^r\{q_j,2q_j,\ldots,Q\},\qquad
 d_\rho=\#\{j:\rho/q_j\in\mathbb Z\}.
\end{equation}
Then $G_{\bq}(p+Q)=\sigma G_{\bq}(p)$, and every positive pole is a
translate of a member of $\PP$. Define its finite Laurent data by
\begin{equation}\label{clc:eq:Laurent}
 G_{\bq}(\rho+z)=\sum_{k=1}^{d_\rho}a_{\rho,k}z^{-k}+O(1).
\end{equation}

\begin{lemma}[Elementary commensurate kernel]\label{clc:lem:kernel}
For $\mu<0$,
\begin{equation}\label{clc:eq:kernel}
 K_{\bq}(\mu)=
 -\frac{\displaystyle\sum_{\rho\in\PP}\sum_{k=1}^{d_\rho}
 a_{\rho,k}e^{\rho\mu}\frac{\mu^{k-1}}{(k-1)!}}
 {1-\sigma e^{Q\mu}}.
\end{equation}
If $\sigma=1$, the simple Laurent coefficients satisfy
\begin{equation}\label{clc:eq:residue-sum}
 \sum_{\rho\in\PP}a_{\rho,1}=0.
\end{equation}
Formula \eqref{clc:eq:kernel} continues to all real $\mu$, with its
removable value at zero understood when $\sigma=1$.
\end{lemma}
\begin{proof}
Take the inverse transform of $G_{\bq}$ on $0<c<q_*$. For $\mu<0$ close
the contour to the right. Choose the closing vertical lines at a fixed
nonpole residue class modulo $Q$. The trigonometric denominators are
then uniformly separated from zero near the real axis, and give
exponential decay at large imaginary part. The closing vertical
integral tends to zero by $e^{\mu\Rea p}$; the horizontal integrals
vanish as their heights tend to infinity. The closure is clockwise,
so the sign of the residue sum is negative.

At $p=\rho+nQ$, the coefficient of $(p-\rho-nQ)^{-k}$ is
$\sigma^n a_{\rho,k}$. The residue of its product with $e^{p\mu}$ is
$\sigma^na_{\rho,k}e^{(\rho+nQ)\mu}\mu^{k-1}/(k-1)!$.
The convergent geometric sum over $n\ge0$ proves
\eqref{clc:eq:kernel}.

When $\sigma=1$, integrate $G_{\bq}$ around a rectangle whose two
vertical edges differ by $Q$, enclosing one positive period of poles.
The vertical integrals cancel by periodicity, and the horizontal ones
vanish by exponential decay. The residue theorem gives
\eqref{clc:eq:residue-sum}. In \eqref{clc:eq:kernel} the numerator
therefore vanishes at zero, so its possible real singularity is
removable. The original inverse transform is holomorphic in a strip
around the real axis. Analytic continuation from $\mu<0$ along the
real axis proves the last assertion. Possible more distant complex
zeros of the denominator are not being declared removable here.
\end{proof}

\begin{theorem}[Finite Lerch closure at every complex order]
\label{clc:thm:finite}
Let $\Rea A>-q_*$ and $b_\rho=(A+\rho)/Q$. For $\mu<0$ and $W\in\C$,
\begin{align}
 \CC_{\bq,A}(W;\mu)
 =-\sum_{\rho\in\PP}\sum_{k=1}^{d_\rho}a_{\rho,k}e^{\rho\mu}
 \sum_{j=0}^{k-1}
 &\frac{(-1)^j\mu^{k-1-j}(W)_j}
 {j!(k-1-j)!Q^{W+j}}\notag\\[-2pt]
 &\hspace{-24mm}\cdot\Phi(\sigma e^{Q\mu},W+j,b_\rho).
 \label{clc:eq:finite}
\end{align}
The whole expression continues to the real axis. If $\sigma=1$ and
$\mu>0$, its two Lerch boundary values agree. At $\mu=0$ the following
central expression is interpreted in its pole-free form below:
\begin{equation}\label{clc:eq:central-raw}
 \CC_{\bq,A}(W;0)=
 \sum_{\rho\in\PP}\sum_{k=1}^{d_\rho}
 \frac{(-1)^ka_{\rho,k}}{(k-1)!Q^{W+k-1}}
 (W)_{k-1}\Phi(\sigma,W+k-1,b_\rho).
\end{equation}
\end{theorem}
\begin{proof}
For $\Rea W>0$, insert \eqref{clc:eq:kernel} into
\eqref{clc:eq:base-fractional}. Expand
$(\mu-u)^{k-1}/(k-1)!$ in powers of $u$, and use
\[
 \frac1{1-\sigma e^{Q(\mu-u)}}
 =\sum_{n=0}^\infty\sigma^ne^{nQ\mu-nQu}.
\]
The absolute convergence is immediate for $\mu<0$ from
$\Rea(A+\rho)>0$. Each Laplace integral is
$\Gamma(W+j)/(A+\rho+nQ)^{W+j}$.
Dividing by $\Gamma(W)$ produces $(W)_j$, proving
\eqref{clc:eq:finite} in the initial half-plane. Since the Lerch argument
has modulus less than one, the right side is entire in $W$. Theorem
\ref{clc:thm:addition} proves equality everywhere.

Initially with $\Rea W$ sufficiently large, taking $\mu\to0^-$ is
justified in the summed integrals and leaves only $j=k-1$. This gives
\eqref{clc:eq:central-raw}. Its extension to all orders follows from
the pole-free identities in Corollary~\ref{clc:cor:polefree}.

For the real continuation when $\sigma=1$, continue the equality from
negative $\mu$ through the upper and lower parts of the holomorphic
strip. The Lerch jump at
$z=e^{Q\mu}>1$ is, up to the common boundary-orientation sign,
\[
 \frac{2\pi i}{\Gamma(W+j)}
 e^{-b_\rho Q\mu}(Q\mu)^{W+j-1}.
\]
For $\Rea(W+j)>0$, this jump follows by indenting the Lerch Euler
integral at its simple pole $u=Q\mu$; its residue is precisely the
displayed coefficient. The general order follows by continuation.
After multiplication by the factors in \eqref{clc:eq:finite}, the
$\rho$ and $j$ dependence in the exponential and in $Q$ cancels.
For fixed $k\ge2$, the remaining coefficient is proportional to
\[
 \sum_{j=0}^{k-1}\frac{(-1)^j}{j!(k-1-j)!}=0.
\]
For $k=1$, the total remaining coefficient is proportional to
$\sum_\rho a_{\rho,1}=0$. This initially proves zero jump away from
the Gamma exceptional points and then everywhere by order continuation.
Thus both boundary values coincide with the analytic ordinary
convolution. For $\sigma=-1$ there is no positive-real Lerch cut to cross.
\end{proof}

\begin{corollary}[Entire central normalization]\label{clc:cor:polefree}
If $\sigma=-1$, every summand of \eqref{clc:eq:central-raw} is entire,
with $\Phi(-1,v,b)=\eta(v,b)$. If $\sigma=1$, choose any fixed reference
$b_*$ with positive real part. Replace the combined $k=1$ terms by
\begin{equation}\label{clc:eq:simple-polefree}
 -Q^{-W}\sum_{\rho\in\PP}a_{\rho,1}
 \bigl[\zeta(W,b_\rho)-\zeta(W,b_*)\bigr].
\end{equation}
For every $k\ge2$ make the replacement
\begin{equation}\label{clc:eq:high-polefree}
 (W)_{k-1}\zeta(W+k-1,b_\rho)
 =(W)_{k-2}\EE_{b_\rho}(W+k-2).
\end{equation}
All these new summands are entire in $W$ and their sum is the desired
central convolution.
\end{corollary}
\begin{proof}
The residues of the two Hurwitz functions in a difference cancel. At
$W=1$ its value is $\psi(b_*)-\psi(b_\rho)$ by
\eqref{clc:eq:stieltjes-convention}. Equation
\eqref{clc:eq:residue-sum} makes subtraction of the reference function
legitimate. For $k\ge2$, the last factor of $(W)_{k-1}$ is
$W+k-2$, giving \eqref{clc:eq:high-polefree}. The function $\EE$ is
entire by its removable normalization. These observations prove the
claim, including every nonpositive-integer resonance.
\end{proof}

The number of terms is finite independently of the orders. It depends
only on the scales and pole multiplicities. No search for an integer
relation among numerical special values enters the proof.

% END sections/03-pole-lattice.tex
% BEGIN sections/04-coefficients.tex
\section{Exact Gamma and Bell formulas for the Laurent coefficients}
\label{clc:sec:coefficients}

The finite closure becomes an explicit algorithm once the coefficients
$a_{\rho,k}$ are known. For a fixed pole put
\[
 I_\rho=\{j:\rho/q_j\in\mathbb Z\},\qquad
 J_\rho=\{1,\ldots,r\}\setminus I_\rho,
 \qquad d=|I_\rho|.
\]
Define the complete exponential Bell polynomials by
\begin{equation}\label{clc:eq:Bell-def}
 \exp\left(\sum_{n\ge1}u_n\frac{z^n}{n!}\right)
 =\sum_{n\ge0}\mathcal B_n(u_1,\ldots,u_n)\frac{z^n}{n!}.
\end{equation}
Thus $\mathcal B_0=1$, $\mathcal B_1=u_1$,
$\mathcal B_2=u_1^2+u_2$, and
$\mathcal B_3=u_1^3+3u_1u_2+u_3$.

\begin{theorem}[Finite coefficient construction]\label{clc:thm:coeff}
The leading coefficient at $\rho$ is
\begin{equation}\label{clc:eq:leading}
 a_{\rho,d}=(-1)^{\sum_{j\in I_\rho}\rho/q_j}
 \prod_{j\in J_\rho}\frac{\pi}{q_j}\csc\frac{\pi\rho}{q_j}.
\end{equation}
For $n\ge1$ let
\begin{align}
 u_n={}&\boldsymbol1_{2\mid n}\,2(n-1)!\zeta(n)
              \sum_{j\in I_\rho}q_j^{-n}\notag\\
 &+\sum_{j\in J_\rho}q_j^{-n}
 \left\{\psi^{(n-1)}(\rho/q_j)
       +(-1)^n\psi^{(n-1)}(1-\rho/q_j)\right\},
 \label{clc:eq:logjets}
\end{align}
where the first line is zero for odd $n$, so it never calls for
$\zeta(1)$. Then, for $1\le k\le d$,
\begin{equation}\label{clc:eq:Bell-coeff}
 a_{\rho,k}=\frac{a_{\rho,d}}{(d-k)!}
                  \mathcal B_{d-k}(u_1,\ldots,u_{d-k}).
\end{equation}
Only $u_1,\ldots,u_{d-1}$ are needed. For rational scales,
$a_{\rho,k}/\pi^{r-k}$ is a real cyclotomic algebraic number.
\end{theorem}
\begin{proof}
Each polar factor has leading term
$(-1)^{\rho/q_j}/z$, proving \eqref{clc:eq:leading}.
The normalized local germ
$z^dG_{\bq}(\rho+z)/a_{\rho,d}$ is analytic and equals one at zero;
its local logarithm is consequently unambiguous.
For a polar factor its logarithm is
\[
 \log\frac{\pi z/q_j}{\sin(\pi z/q_j)}
 =\sum_{h=1}^\infty\frac{\zeta(2h)}{h}\left(\frac z{q_j}\right)^{2h}.
\]
This follows from the sine product, or from Gamma reflection and the
local Gamma series. Its $n$th derivative at zero is the first line of
\eqref{clc:eq:logjets}.
For a nonpolar factor use
\[
 g_{q_j}(\rho+z)=q_j^{-1}
 \Gamma\left(\frac{\rho+z}{q_j}\right)
 \Gamma\left(1-\frac{\rho+z}{q_j}\right).
\]
Differentiating its local logarithm gives the second line. The arguments
$1-\rho/q_j$ here may be negative, but are not integers; the meromorphic
polygamma values are well defined. No global branch of their separate
Gamma logarithms is required. Exponentiating the normalized local
logarithm gives \eqref{clc:eq:Bell-coeff}.

For the final assertion one can avoid transcendental cancellation
entirely. Expand the sine factors directly at rational multiples of
$\pi$, invert their finite Taylor polynomials, and multiply to order
$d-1$. The trigonometric constants are real cyclotomic algebraic
numbers, and the Taylor powers supply exactly $\pi^{r-k}$.
\end{proof}

The two constructions in this proof are independent computational
routes. The supplied code uses the logarithmic Bell recurrence in
\texttt{pole\_table}, and direct reciprocal-sine Taylor multiplication in
\texttt{direct\_local\_table}. Exact comparisons are made after symbolic
simplification, not after floating-point evaluation of polygamma terms.

\subsection{A genuinely mixed three-scale example}
For $\bq=(1,2,3)$, a common period is $Q=6$ and $\sigma=-1$.
The complete table is as follows; blank entries mean zero or absent
higher pole order.
\begin{center}
\begin{tabular}{c c c c}
\toprule
$\rho$ & $a_{\rho,1}$ & $a_{\rho,2}$ & $a_{\rho,3}$\\
\midrule
$1$ & $-\sqrt3\pi^2/9$ & &\\
$2$ & $-2\pi^2/27$ & $-2\sqrt3\pi/9$ &\\
$3$ & $0$ & $-\pi/2$ &\\
$4$ & $2\pi^2/27$ & $-2\sqrt3\pi/9$ &\\
$5$ & $\sqrt3\pi^2/9$ & &\\
$6$ & $-49\pi^2/216$ & $0$ & $-1$\\
\bottomrule
\end{tabular}
\end{center}
Together with \eqref{clc:eq:central-raw}, this is an explicit finite
identity for every total order and every $\Rea A>-1$, using alternating
Hurwitz functions at $(A+\rho)/6$. No additional infinite sum is hidden
in the coefficient construction.
At $W=0$ every $k\ge2$ term vanishes because $(0)_{k-1}=0$, and
$\eta(0,b)=1/2$. Thus
\begin{equation}\label{clc:eq:123zero}
 K_{(1,2,3)}(0)=-\frac12\sum_\rho a_{\rho,1}
 =\frac{49\pi^2}{432}.
\end{equation}
The independent parity theorem in Section~\ref{clc:sec:parity} gives
exactly the same value without constructing this table.

The coefficients also obey a useful scaling check. If all scales and
poles are multiplied by $\lambda>0$, then
\begin{equation}\label{clc:eq:coef-scaling}
 a^{\lambda\bq}_{\lambda\rho,k}
 =\lambda^{k-r}a^{\bq}_{\rho,k}.
\end{equation}
This follows at once from
$G_{\lambda\bq}(p)=\lambda^{-r}G_{\bq}(p/\lambda)$ and is included
in the exact verification suite.

% END sections/04-coefficients.tex
% BEGIN sections/05-two-scales.tex
\section{The complete integer two-scale identity}
\label{clc:sec:two}

The general table construction has a particularly short closed form
for scales $1$ and $m$, where $m$ is any positive integer. Write
\begin{equation}\label{clc:eq:pair-data}
 \sigma_m=(-1)^{m+1},\qquad
 c_r^{(m)}=(-1)^{r+1}\csc\frac{\pi r}{m}\quad(1\le r<m).
\end{equation}
When $m=1$ all sums over $r$ below are empty.
For $\Rea A>-1$, define the \emph{unnormalized} integral
\begin{equation}\label{clc:eq:I-def}
 I_m(s,t;A)=\int_0^\infty
 F_{1+A,s}(y)F_{1+A/m,t}(y^{-m})\frac{\dd y}{y}.
\end{equation}
The normalized quantity is
$D_m(W;A)=\CC_{(1,m),A}(W;0)=m^{-t}I_m(s,t;A)$.

\begin{theorem}[All integer scale ratios and all complex orders]
\label{clc:thm:pair}
For $s,t\in\C$, $W=s+t$, and $\Rea A>-1$,
\begin{align}
 I_m(s,t;A)=m^{-s-1}\biggl[&
 \pi\sum_{r=1}^{m-1}c_r^{(m)}
       \Phi\left(\sigma_m,W,\frac{A+r}{m}\right)\notag\\
 &+\sigma_m W\Phi\left(\sigma_m,W+1,1+\frac A m\right)\biggr].
 \label{clc:eq:pair}
\end{align}
The integral is ordinary, absolutely convergent, and entire in $s,t$.
For odd $m$, the last term is $\EE_{1+A/m}(W)$, and the simple Hurwitz
sum is grouped into differences to remove its pole at $W=1$.
\end{theorem}
\begin{proof}
Use $Q=m$. The poles at $r=1,\ldots,m-1$ are simple, with
\[
 a_{r,1}=(-1)^r\frac\pi m\csc\frac{\pi r}{m}.
\]
At $\rho=m$ both factors are polar, and
$a_{m,2}=(-1)^{m+1}=\sigma_m$, $a_{m,1}=0$.
Substitution in \eqref{clc:eq:central-raw} gives the normalized formula
with prefactor $m^{-W-1}$. Multiplication by $m^t$ gives
\eqref{clc:eq:pair}. Entire dependence and absolute convergence were
proved before the residue calculation. For odd $m$,
$c_{m-r}^{(m)}=-c_r^{(m)}$, so $\sum_rc_r^{(m)}=0$ and the
simple-pole cancellation is explicit.
\end{proof}

\subsection{Scales two and three}
At $A=0$ the functions in the integrand are ordinary polylogarithms.
For $m=2$, the two reductions
$\eta(W,1/2)=2^W\beta(W)$ and $\eta(W,1)=\eta(W)$ give
\eqref{clc:eq:intro-beta}. For $m=3$, let $\chi_{-3}$ be the real
nonprincipal character modulo three and set
\[
 L(W,\chi_{-3})=3^{-W}
       [\zeta(W,1/3)-\zeta(W,2/3)].
\]
Then
\begin{equation}\label{clc:eq:pair3}
 \boxed{\displaystyle
 \int_0^\infty\Li_s(-y)\Li_t(-y^{-3})\frac{\dd y}{y}
 =\frac{2\pi\,3^{t-1}}{\sqrt3}L(W,\chi_{-3})
    +3^{-s-1}\EE_1(W).}
\end{equation}
For $s=t=1$ this yields
\begin{equation}\label{clc:eq:log3}
 \int_0^\infty\log(1+y)\log(1+y^{-3})\frac{\dd y}{y}
 =\frac{2\pi}{\sqrt3}L(2,\chi_{-3})+\frac29\zeta(3).
\end{equation}
The first two numerical values are
\begin{align*}
 I_2(1,1;0)&=2.426819442923313294752678355954754053578\ldots,\\
 I_3(1,1;0)&=3.101375395830302048787723234392638\ldots.
\end{align*}
They were checked against direct real-line logarithmic quadrature;
they are diagnostics, not the proof.

\subsection{Zero-order resonance and an explicit normalization obstruction}
At zero orders one obtains
\begin{equation}\label{clc:eq:zero-examples}
 I_2(0,0;0)=\frac\pi4,\qquad
 I_3(0,0;0)=\frac13+\frac{2\pi}{9\sqrt3}.
\end{equation}
The $1/3$ in the second expression is the value
$3^{-1}\EE_1(0)$. Substituting $W=0$ in a separated factor
$W\zeta(W+1)$ and declaring it zero would lose this term.

Even at the smallest nonzero total order the raw unequal-scale
integral fails pure order addition:
\begin{align}
 I_2(1,0;0)&=\frac{5\pi^2}{48},&
 I_2(0,1;0)&=\frac{5\pi^2}{24}.
 \label{clc:eq:swap}
\end{align}
For instance, use $\beta(1)=\pi/4$ and $\eta(2)=\pi^2/12$ in
\eqref{clc:eq:intro-beta}. These two values have the same total order
but differ by a factor of two. The normalized integrals agree, as
Theorem~\ref{clc:thm:addition} requires. This is an exact counterexample
to a tempting extension of the old equal-scale statement, not an error
attributed to that statement itself.

For $m=1$, Theorem~\ref{clc:thm:pair} reduces to the preceding report's
identity $I_1(s,t;A)=\EE_{1+A}(s+t)$. This limiting consistency is useful,
but the equal-scale case is prior work, not a new claim.

% END sections/05-two-scales.tex
% BEGIN sections/06-harmonic.tex
\section{Balanced harmonic remainders and their order derivatives}
\label{clc:sec:harmonic}

The scale alignment also identifies the correct unequal truncation
lengths. For $N\ge0$ put
\begin{equation}\label{clc:eq:R}
 \mathcal R_{N,s}(y)=\Li_s(-y)-\sum_{k=1}^{N}\frac{(-y)^k}{k^s}.
\end{equation}
The empty sum is zero. For $a>0$ and $\sigma\in\{1,-1\}$ write
\begin{equation}\label{clc:eq:Hcolored}
 H_{N,\sigma}^{(v)}(a)=\sum_{k=0}^{N-1}\frac{\sigma^k}{(k+a)^v}.
\end{equation}
This finite function is entire in $v$. At $a=1$, $\sigma=1$, it is the
usual generalized harmonic number $H_N^{(v)}$.

\begin{theorem}[Unequal-scale harmonic-remainder product]
\label{clc:thm:remainder}
For all integers $m\ge1$, $N\ge0$, and all complex $s,t$, with $W=s+t$,
\begin{align}
 &\int_0^\infty\mathcal R_{mN,s}(y)
            \mathcal R_{N,t}(y^{-m})\frac{\dd y}{y}
 = I_m(s,t;0)\notag\\
 &\quad-m^{-s-1}\biggl[
 \pi\sum_{r=1}^{m-1}c_r^{(m)}
 H_{N,\sigma_m}^{(W)}\left(\frac r m\right)
 +\sigma_m W H_{N,\sigma_m}^{(W+1)}(1)\biggr].
 \label{clc:eq:remainder}
\end{align}
The integral is absolutely convergent and entire in its two orders.
Every independent order derivative can be passed through it.
\end{theorem}
\begin{proof}
The elementary shift of the Lerch series, first for $0<y<1$ and then
by continuation along $y>0$, gives
\begin{equation}\label{clc:eq:Rshift}
 F_{N+1,s}(y)=-(-y)^{-N}\mathcal R_{N,s}(y).
\end{equation}
Take $A=mN$ in \eqref{clc:eq:I-def}. The two factors in
\eqref{clc:eq:Rshift} have powers $y^{-mN}$ and $y^{mN}$, so they
cancel; their signs give
\[
 \int_0^\infty\mathcal R_{mN,s}(y)
       \mathcal R_{N,t}(y^{-m})\frac{\dd y}{y}
 =\sigma_m^N I_m(s,t;mN).
\]
This also proves convergence and all order-holomorphic assertions by
Theorem~\ref{clc:thm:pair} and the weighted lemma.

The finite shift formula is
\[
 \Phi(\sigma,v,a+N)=\sigma^N
             [\Phi(\sigma,v,a)-H_{N,\sigma}^{(v)}(a)]
 \quad(\sigma=\pm1).
\]
It follows directly by removing the first $N$ terms; at $\sigma=1$
it is continued through the usual Hurwitz pole in the combined formula.
All arguments in \eqref{clc:eq:pair} increase by $N$ when $A=mN$.
The resulting factor $\sigma_m^N$ cancels the preceding sign, proving
\eqref{clc:eq:remainder}. All harmonic corrections are finite entire
functions of the orders, so no additional continuation issue remains.
\end{proof}

An explicit example is
\begin{align}
 &\int_0^\infty
 \left[-\log(1+y)+y-\frac{y^2}{2}\right]
 \left[-\log(1+y^{-2})+y^{-2}\right]\frac{\dd y}{y}\notag\\
 &\hspace{25mm}=\pi G-\frac38\zeta(3)-\pi+\frac12.
 \label{clc:eq:Rexample}
\end{align}
This is the case $m=2$, $N=1$, $s=t=1$. The two brackets must be kept
together. Expanding the product and integrating its monomials separately
would introduce divergent integrals and would not prove the identity.
The unequal truncation lengths $2$ and $1$ are dictated by the scales;
they are not interchangeable with two equal truncation lengths.

\subsection{Logarithmic harmonic derivatives}
For every $h\ge0$,
\begin{equation}\label{clc:eq:Hjets}
 \partial_v^hH_{N,\sigma}^{(v)}(a)
 =(-1)^h\sum_{k=0}^{N-1}
       \frac{\sigma^k\log^h(k+a)}{(k+a)^v}.
\end{equation}
Consequently the derivatives of \eqref{clc:eq:remainder} give explicit
identities for products of polylogarithmic order derivatives and finite
harmonic-logarithmic corrections, with no new infinite coefficients.
The prefactor $m^{-s-1}$ contributes powers of $-\log m$ and must be
differentiated as well.

For example, the first $s$ derivative of the correction bracket in
\eqref{clc:eq:remainder}, evaluated at any pair $(s,t)$, is
\begin{align*}
 m^{-s-1}\biggl\{&\pi\sum_r c_r^{(m)}
       [\partial_W H_{N,\sigma_m}^{(W)}(r/m)
                     -\log m\,H_{N,\sigma_m}^{(W)}(r/m)]\\
 &+\sigma_m[(1-W\log m)H_{N,\sigma_m}^{(W+1)}(1)
                  +W\partial_W H_{N,\sigma_m}^{(W+1)}(1)]\biggr\}.
\end{align*}
The $t$ derivative has no derivative of the prefactor. Thus the explicit
harmonic identities retain the same scale asymmetry as the raw
polylogarithm integral. Arbitrarily high derivatives follow from the
finite Leibniz rule rather than an unproved extrapolation.

% END sections/06-harmonic.tex
% BEGIN sections/07-moments.tex
\section{Finite closure of every mixed logarithmic moment}
\label{clc:sec:moments}

The total-order theorem concerns the unweighted convolution. Polynomial
weights in the logarithmic variables still have finite closure, even
though they distinguish the independent orders.
For a multi-index $\boldsymbol\ell=(\ell_1,\ldots,\ell_r)$ of
nonnegative integers, put $L=\sum_j\ell_j$ and define
\begin{equation}\label{clc:eq:M}
 \mathcal M_{\boldsymbol\ell}(\bs;A;\mu)=
 \int_{\R^{r-1}}\prod_{j=1}^r
      x_j^{\ell_j}h_{q_j,A,s_j}(x_j)\dd\boldsymbol x,
 \qquad x_r=\mu-\sum_{j<r}x_j.
\end{equation}
The measure denotes the first $r-1$ variables as before.

\begin{theorem}[A terminating mixed-moment algorithm]
\label{clc:thm:moments}
Let the scales be commensurate and $\Rea A>-q_*$. Then
\eqref{clc:eq:M} is absolutely convergent and entire in all orders.
For $0\le D\le L$, form the meromorphic trigonometric function
\begin{equation}\label{clc:eq:GD}
 G_D(p;\bs)=
 \sum_{\substack{0\le d_j\le\ell_j\\\sum_jd_j=D}}
 \prod_{j=1}^r
 \binom{\ell_j}{d_j}(s_j)_{d_j}
 (-\partial_p)^{\ell_j-d_j}g_{q_j}(p).
\end{equation}
Compute the finite Laurent table of each $G_D$ on the same period $Q$.
Apply formula \eqref{clc:eq:finite}, or its pole-free central version,
to that table with order $W+D$. Summing the answers over $D$ gives
exactly $\mathcal M_{\boldsymbol\ell}$.
The pole set remains $\PP$, and no pole order exceeds $r+L-D$.
Every fixed mixed order derivative of this finite formula is valid.
\end{theorem}
\begin{proof}
Multiplication by $x^{\ell_j}$ corresponds to
$(-\partial_p)^{\ell_j}$ on the bilateral transform. Lemma
\ref{clc:lem:weighted} justifies these operations and the integrals.
The Leibniz rule gives
\[
 (-\partial_p)^{\ell_j}[g_{q_j}(p)(p+A)^{-s_j}]
 =\sum_{d=0}^{\ell_j}\binom{\ell_j}{d}(s_j)_d
   (-\partial_p)^{\ell_j-d}g_{q_j}(p)(p+A)^{-s_j-d}.
\]
Multiplying and grouping by $D$ yields
$\sum_DG_D(p;\bs)(p+A)^{-W-D}$.
Every derivative of a cosecant has the same signed period as the
cosecant, so $G_D(p+Q;\bs)=\sigma G_D(p;\bs)$.
Differentiation creates no new poles and increases pole orders by at
most the number of derivatives. The total is at most $r+L-D$.
All these functions retain exponential vertical decay.

The residue proof of Lemma~\ref{clc:lem:kernel} and Theorem
\ref{clc:thm:finite} uses precisely these properties, not a special
factorization into undifferentiated cosecants. Thus it applies to each
$G_D$. In the periodic case the simple residues still sum to zero by
the period rectangle. The same pole-removal and branch-jump proofs
therefore remain valid. The resulting expression is a finite sum of
entire functions of the orders with polynomial coefficient factors.
The weighted holomorphy proves that it also equals every differentiated
ordinary integral claimed in the statement.
\end{proof}

\subsection{A convenient differential-polynomial implementation}
For a single scale set $\kappa_q=(\pi/q)^2$ and define
\begin{align}
 R_0(s;T,U;\kappa)&=1,\notag\\
 R_{\ell+1}&=(T+sU)R_\ell
       +(\kappa+T^2)\partial_TR_\ell+U^2\partial_UR_\ell.
 \label{clc:eq:Rrec}
\end{align}
With $T=(\pi/q)\cot(\pi p/q)$ and $U=(p+A)^{-1}$ one has
\begin{equation}\label{clc:eq:Ridentity}
 (-\partial_p)^\ell[g_q(p)(p+A)^{-s}]
 =g_q(p)(p+A)^{-s}R_\ell(s;T,U;\kappa_q).
\end{equation}
Indeed $g_q'=-Tg_q$, $T'=-(\kappa_q+T^2)$ and $U'=-U^2$;
the product rule proves the recurrence. The first cases are
\begin{align*}
 R_1&=T+sU,\\
 R_2&=2T^2+2sTU+(s)_2U^2+\kappa,\\
 R_3&=6T^3+6sT^2U+3(s)_2TU^2+(s)_3U^3
                  +5\kappa T+3s\kappa U.
\end{align*}
Different scales have different $T$ and $\kappa$ values. Setting every
$\kappa$ to $\pi^2$ would silently revert to equal scales.
Formula \eqref{clc:eq:GD} avoids that error and provides an unambiguous
finite normal form before pole extraction.

There are also exact consistency constraints. At every $\mu$,
\begin{equation}\label{clc:eq:moment-sum}
 \sum_{j=1}^r\mathcal M_{\boldsymbol e_j}(\bs;A;\mu)
       =\mu\CC_{\bq,A}(W;\mu).
\end{equation}
This follows directly from $\sum_jx_j=\mu$ in the ordinary integral.
Unlike the equal-scale first-moment formulas, the individual terms need
not be a fixed $1/r$ share. The identity is a check on their sum, not a
claim that scale asymmetry disappears.

% END sections/07-moments.tex
% BEGIN sections/08-stieltjes.tex
\section{All-index Stieltjes identities and Gamma specializations}
\label{clc:sec:stieltjes}

\subsection{Ordinary products of normalized order jets}
Define the normalized order jets
\begin{equation}\label{clc:eq:Ljet}
 \mathscr L_{q,n}^{A}(x)=
       \left.\partial_s^n h_{q,A,s}(x)\right|_{s=0}.
\end{equation}
The normalization $q^{-s}$ is part of this definition; its derivatives
contribute powers of $-\log q$. For a multi-index
$\boldsymbol n=(n_1,\ldots,n_r)$, $M=\sum_jn_j$, Theorem
\ref{clc:thm:addition} gives the ordinary product law
\begin{equation}\label{clc:eq:jet-product}
 \int_{\R^{r-1}}\prod_{j=1}^r\mathscr L_{q_j,n_j}^{A}(x_j)
                  \dd\boldsymbol x
 =\CC_{\bq,A}^{[M]}(0;0),\qquad x_r=-\sum_{j<r}x_j.
\end{equation}
It is absolutely convergent. In particular, even at unequal scales the
answer depends on the distribution of order derivatives only through
their total $M$. It is not a Hermitian product, so no positivity of
these Stieltjes expressions is implied.

Here are the exact rules that turn every central closure into Stieltjes
and zeta derivatives:
\begin{align}
 \EE_a^{[h]}(0)&=
 \begin{cases}1,&h=0,\\(-1)^{h-1}h\gamma_{h-1}(a),&h\ge1,
 \end{cases}\label{clc:eq:Ejet}\\
 \EE_a^{[h]}(v)&=v\zeta^{[h]}(v+1,a)
                +h\zeta^{[h-1]}(v+1,a)\quad(v\ne0).
 \label{clc:eq:Ejet-nonzero}
\end{align}
The second term in \eqref{clc:eq:Ejet-nonzero} is absent when $h=0$.
The first identity is just the Laurent convention
\eqref{clc:eq:stieltjes-convention} multiplied by $v$.

For example, when $\sigma=1$ an explicit all-index coefficient formula
for the right side of \eqref{clc:eq:jet-product} is
\begin{align}
 M![z^M]Q^{-z}\biggl\{&-\sum_\rho a_{\rho,1}
         [\zeta(z,b_\rho)-\zeta(z,b_*)]\notag\\
 &+\sum_\rho\sum_{k=2}^{d_\rho}
 \frac{(-1)^ka_{\rho,k}}{(k-1)!Q^{k-1}}
      (z)_{k-2}\EE_{b_\rho}(z+k-2)\biggr\}.
 \label{clc:eq:alljet-coeff}
\end{align}
This is finite coefficient extraction from explicitly given analytic
germs. The only Stieltjes constants in it arise from
\eqref{clc:eq:Ejet}; the other terms involve ordinary zeta order
derivatives at zero or positive integers. For $\sigma=-1$, use instead
\begin{equation}\label{clc:eq:alljet-eta}
 M![z^M]Q^{-z}\sum_\rho\sum_{k=1}^{d_\rho}
 \frac{(-1)^ka_{\rho,k}}{(k-1)!Q^{k-1}}
             (z)_{k-1}\eta(z+k-1,b_\rho).
\end{equation}
These formulas, together with finite product differentiation, are
explicit at every index; they do not require a recurrence fitted to
low-order data.

\subsection{Generalized harmonic coefficients}
For $j\ge1$ the exact local identity
\begin{equation}\label{clc:eq:poch-harmonic}
 (z)_j=(j-1)!z\exp\left(
       \sum_{h=1}^\infty\frac{(-1)^{h-1}}h H_{j-1}^{(h)}z^h\right)
\end{equation}
follows by taking the logarithm of
$\prod_{k=1}^{j-1}(1+z/k)$. Consequently, for $1\le l\le j$,
\begin{equation}\label{clc:eq:poch-Bell}
 \left.\partial_z^l(z)_j\right|_{z=0}
 =l(j-1)!\mathcal B_{l-1}
 \bigl(H_{j-1},-H_{j-1}^{(2)},2!H_{j-1}^{(3)},\ldots\bigr).
\end{equation}
Derivatives of order greater than $j$ vanish because $(z)_j$ is a
polynomial. Equations \eqref{clc:eq:alljet-coeff}--\eqref{clc:eq:poch-Bell}
therefore express every jet by finite generalized-harmonic polynomials,
zeta derivatives, and Stieltjes constants. At other integer resonances,
the finite polynomial should be shifted and expanded directly; dividing
by a vanishing Pochhammer factor is unnecessary and unsafe.

\subsection{A short all-index formula for two scales}
Put $b_r=(A+r)/m$ and $b_m=1+A/m$. Define
$e_0(b)=1$ and $e_h(b)=(-1)^{h-1}h\gamma_{h-1}(b)$ for $h\ge1$.
For odd $m$, Theorem~\ref{clc:thm:pair} gives
\begin{equation}\label{clc:eq:odd-pair-jet}
 \boxed{\displaystyle
 D_m^{[M]}(0;A)=\frac1m\sum_{h=0}^M\binom Mh(-\log m)^{M-h}
 \left[\pi\sum_{r=1}^{m-1}c_r^{(m)}\zeta^{[h]}(0,b_r)
                       +e_h(b_m)\right].}
\end{equation}
For even $m$, the corresponding formula is
\begin{equation}\label{clc:eq:even-pair-jet}
 \boxed{\displaystyle
 D_m^{[M]}(0;A)=\frac1m\sum_{h=0}^M\binom Mh(-\log m)^{M-h}
 \left[\pi\sum_{r=1}^{m-1}c_r^{(m)}\eta^{[h]}(0,b_r)
               -h\eta^{[h-1]}(1,b_m)\right],}
\end{equation}
where the last term is zero when $h=0$. These identities follow from
the ordinary product rule applied to
$D_m(W;A)=m^{-W-1}[\cdots]$ in \eqref{clc:eq:pair}; all differentiated
integrals are justified by the weighted lemma.

The apparently separate eta jets at one are themselves finite
Stieltjes differences:
\begin{equation}\label{clc:eq:eta-Stieltjes}
 \eta^{[h]}(1,b)=\frac12\sum_{j=0}^h\binom hj(-\log2)^{h-j}(-1)^j
 \left[\gamma_j\left(\frac b2\right)
       -\gamma_j\left(\frac{b+1}2\right)\right].
\end{equation}
Indeed the two Hurwitz poles in \eqref{clc:eq:eta} cancel before
multiplication by $2^{-v}$; expansion at $v=1$ then gives the formula.
For $h=0$ it agrees with
\begin{equation}\label{clc:eq:eta-one}
 \eta(1,b)=\frac12\left[\psi\left(\frac{b+1}2\right)
                              -\psi\left(\frac b2\right)\right].
\end{equation}
This check fixes the sign of every Stieltjes difference.

\subsection{Quarter-Gamma and third-Gamma product identities}
The classical identity
$\zeta^{[1]}(0,b)=\log\Gamma(b)-\tfrac12\log(2\pi)$
\cite{DLMF25} gives
\begin{equation}\label{clc:eq:eta-zerojet}
 \eta^{[1]}(0,b)=\log\Gamma(b/2)-\log\Gamma((b+1)/2)
                              -\frac12\log2.
\end{equation}
Here and below the logarithm of Gamma is the holomorphic one on the
right half-plane, real on the positive axis.
At $A=0$, two explicit consequences are
\begin{align}
 D_2^{[1]}(0;0)
 &=\frac\pi2\log\frac{\Gamma(1/4)}{\Gamma(3/4)}
                  -\frac{\pi+1}{2}\log2,\label{clc:eq:quarter}\\
 D_3^{[1]}(0;0)
 &=\frac{2\pi}{3\sqrt3}\log\frac{\Gamma(1/3)}{\Gamma(2/3)}
   -\frac{2\pi\log3}{9\sqrt3}+\frac{\gamma-\log3}{3}.
 \label{clc:eq:third}
\end{align}
By \eqref{clc:eq:jet-product}, the first value is, for example,
\[
 \int_\R\mathscr L_{1,1}^{0}(x)\mathscr L_{2,0}^{0}(-x)\dd x
 =\int_\R\mathscr L_{1,0}^{0}(x)\mathscr L_{2,1}^{0}(-x)\dd x.
\]
The derivative inside the second integral includes the scale
normalization, so this does not contradict \eqref{clc:eq:swap}.

One order higher, \eqref{clc:eq:pair3} in normalized form gives
\begin{equation}\label{clc:eq:third2}
 D_3^{[2]}(0;0)=\frac{2\pi}{3\sqrt3}L^{[2]}(0,\chi_{-3})
       +\frac{(\log3)^2-2\gamma\log3-2\gamma_1}{3}.
\end{equation}
Thus an ordinary product of first order jets, with two different scales,
contains a Dirichlet $L$-derivative and the first Stieltjes constant in
one exact formula. General $A$ in \eqref{clc:eq:odd-pair-jet} and
\eqref{clc:eq:even-pair-jet} replaces these by generalized Stieltjes
constants and shifted Gamma values.

% END sections/08-stieltjes.tex
% BEGIN sections/09-primitives.tex
\section{Argument derivatives and all anchored primitive orders}
\label{clc:sec:primitives}

Order differentiation is distinct from differentiating the center or
the logarithmic argument. The contour formula makes the latter two
operations equally explicit:
\begin{align}
 \partial_A\CC_{\bq,A}(W;\mu)
      &=-W\CC_{\bq,A}(W+1;\mu),\label{clc:eq:A-ladder}\\
 (\partial_\mu+A)\CC_{\bq,A}(W;\mu)
      &=\CC_{\bq,A}(W-1;\mu).
 \label{clc:eq:mu-ladder}
\end{align}
Both follow by differentiating \eqref{clc:eq:C-contour}; its exponential
vertical majorant justifies each derivative. They also follow from the
corresponding profile identities. Iteration gives
\begin{equation}\label{clc:eq:A-higher}
 \partial_A^d\CC_{\bq,A}(W;\mu)
       =(-1)^d(W)_d\CC_{\bq,A}(W+d;\mu).
\end{equation}
Every mixed spectral derivative can be applied to these identities.
The Pochhammer factor must then be differentiated too.

\subsection{First and higher anchored antiderivatives}
For an anchor $A_0$ in $\Rea A>-q_*$,
\begin{equation}\label{clc:eq:first-primitive}
 \int_{A_0}^A\CC_{\bq,u}(W;\mu)\dd u
 =\frac{\CC_{\bq,A_0}(W-1;\mu)
                  -\CC_{\bq,A}(W-1;\mu)}{W-1}.
\end{equation}
The integral is along any path in that half-plane. The right side is
filled in at $W=1$ by
\begin{equation}\label{clc:eq:first-resonant}
 \CC_{\bq,A_0}^{[1]}(0;\mu)-\CC_{\bq,A}^{[1]}(0;\mu),
\end{equation}
since $\CC_{\bq,A}(0;\mu)=K_{\bq}(\mu)$ is independent of $A$.
At $\mu=0$, this already gives a finite Gamma and polygamma primitive
of the weight-one convolution.

The same mechanism supplies every primitive order. Define
\begin{align}
 T_d(W;A,A_0;\mu)={}&\CC_{\bq,A}(W-d;\mu)\notag\\
 &-\sum_{j=0}^{d-1}\frac{(A-A_0)^j}{j!}(-1)^j(W-d)_j
             \CC_{\bq,A_0}(W-d+j;\mu).
 \label{clc:eq:T-d}
\end{align}

\begin{theorem}[All anchored primitives, including resonance]
\label{clc:thm:primitives}
For $d\ge1$,
\begin{equation}\label{clc:eq:primitive-d}
 \boxed{\displaystyle
 \int_{A_0}^A\frac{(A-u)^{d-1}}{(d-1)!}
                 \CC_{\bq,u}(W;\mu)\dd u
 =\frac{(-1)^dT_d(W;A,A_0;\mu)}{(W-d)_d}.}
\end{equation}
The right side is entire in $W$ after removing its apparent poles at
$W=1,\ldots,d$. At $W=k$ in that range its value is
\begin{equation}\label{clc:eq:primitive-resonance}
 \frac{(-1)^k}{(k-1)!(d-k)!}
            \left.\partial_WT_d(W;A,A_0;\mu)\right|_{W=k}.
\end{equation}
All order derivatives of these primitive identities are valid.
\end{theorem}
\begin{proof}
Apply the ordinary Taylor formula with integral remainder, in the
variable $A$, to $\CC_{\bq,A}(W-d;\mu)$ at $A_0$. Its $j$th derivative
is given by \eqref{clc:eq:A-higher}. Its $d$th derivative is
$(-1)^d(W-d)_d\CC_{\bq,A}(W;\mu)$. This proves
\eqref{clc:eq:primitive-d} when the polynomial denominator is nonzero.
The left side is entire in $W$, by local uniform holomorphy on a compact
integration path. Therefore all apparent poles are removable.
The polynomial $(W-d)_d=\prod_{j=1}^d(W-j)$ has simple zeros, with
its derivative at $W=k$ equal to
$(-1)^{d-k}(k-1)!(d-k)!$.
Taking the quotient limit gives \eqref{clc:eq:primitive-resonance}.
Local uniform holomorphy and Cauchy's formula give every order derivative.
\end{proof}

This formulation fixes all integration constants by the anchor and
its vanishing derivatives through order $d-1$. It is stronger than
writing an unanchored formal primitive and leaving its constants
ambiguous at resonant orders.

\subsection{Explicit Gamma primitives of the two-scale family}
For $D_m(W;A)=\CC_{(1,m),A}(W;0)$, define $b_r=(A+r)/m$.
When $m$ is odd, put
\begin{equation}\label{clc:eq:Podd}
 P_m(A)=\frac1m\left[\psi(b_m)
                  -\pi\sum_{r=1}^{m-1}c_r^{(m)}\log\Gamma(b_r)\right].
\end{equation}
When $m$ is even, put
\begin{equation}\label{clc:eq:Peven}
 P_m(A)=\frac1m\left[\eta(1,b_m)-\pi\sum_{r=1}^{m-1}c_r^{(m)}
 \log\frac{\Gamma(b_r/2)}{\Gamma((b_r+1)/2)}\right].
\end{equation}
Then in both cases
\begin{equation}\label{clc:eq:Pderivative}
 P_m'(A)=D_m(1;A),\qquad
 \int_{A_0}^AD_m(1;u)\dd u=P_m(A)-P_m(A_0).
\end{equation}
To check the odd case directly, use
$\psi'(b)=\zeta(2,b)$ and
$\sum_rc_r^{(m)}\zeta(1,b_r)=-\sum_rc_r^{(m)}\psi(b_r)$,
where the sum is understood after its residue cancellation.
For the even case use
$\partial_b\eta(1,b)=-\eta(2,b)$ and
\[
 \partial_b\log\frac{\Gamma(b/2)}{\Gamma((b+1)/2)}=-\eta(1,b).
\]
Both derivatives reproduce the $W=1$ case of the normalized pair formula,
including its $m^{-2}$ factor. This is also an independent check of the
sign in \eqref{clc:eq:first-resonant}.

Subsequent center derivatives of \eqref{clc:eq:Podd} and
\eqref{clc:eq:Peven} are finite polygamma expressions. The all-index
Stieltjes formulas in Section~\ref{clc:sec:stieltjes}, followed by
\eqref{clc:eq:primitive-d}, give their anchored order-jet primitives.
Thus differentiation and integration are both covered without replacing
a pole-cancelled expression by its singular separate pieces.

% END sections/09-primitives.tex
% BEGIN sections/10-unequal-centers.tex
\section{Unequal centers: normal convergence through every order jet}
\label{clc:sec:unequal}

The preceding finite identities share one affine center $A$.
For independent centers $A_j$, define
\begin{equation}\label{clc:eq:U}
 \UU_{\bq,\bA}(\bs;\mu)=
 \int_{\R^{r-1}}\prod_{j=1}^r h_{q_j,A_j,s_j}(x_j)
             \dd\boldsymbol x.
\end{equation}
We require a common weight:
\begin{equation}\label{clc:eq:unequal-strip}
 \max\bigl(0,-\min_j\Rea A_j\bigr)<c<q_*.
\end{equation}
This condition is stronger than the separate requirements
$\Rea A_j>-q_j$ when the scales are unequal; it must not be dropped.
It holds, in particular, if all $\Rea A_j>-q_*$.
The weighted proof of Theorem~\ref{clc:thm:addition} applies to give
\begin{equation}\label{clc:eq:U-contour}
 \UU_{\bq,\bA}(\bs;\mu)=\frac1{2\pi i}
 \int_{c-i\infty}^{c+i\infty}e^{p\mu}G_{\bq}(p)
                \prod_{j=1}^r(p+A_j)^{-s_j}\dd p.
\end{equation}
It is entire in its independent orders; no integer-order restriction is
being analytically differentiated.

Choose a real reference $A$ with $A+c>0$ and set $\delta_j=A_j-A$.
Suppose
\begin{equation}\label{clc:eq:rho-condition}
 \rho_0=\max_j\frac{|\delta_j|}{A+c}<1.
\end{equation}
Define explicit polynomial coefficients by
\begin{align}
 \prod_{j=1}^r(1+\delta_jz)^{-s_j}
   &=\sum_{N=0}^\infty h_N(\bs;\bdelta)z^N,\label{clc:eq:hN-gen}\\
 h_N(\bs;\bdelta)
   &=\sum_{n_1+\cdots+n_r=N}
       \prod_{j=1}^r\frac{(-\delta_j)^{n_j}(s_j)_{n_j}}{n_j!}.
 \label{clc:eq:hN}
\end{align}
In particular $h_0=1$, $h_1=-\sum_js_j\delta_j$, and
\[
 h_2=\frac12\sum_js_j(s_j+1)\delta_j^2
                 +\sum_{i<j}s_is_j\delta_i\delta_j.
\]

\begin{theorem}[Normally convergent unequal-center closure]
\label{clc:thm:unequal}
Under \eqref{clc:eq:unequal-strip} and \eqref{clc:eq:rho-condition},
\begin{equation}\label{clc:eq:unequal-series}
 \boxed{\displaystyle
 \UU_{\bq,\bA}(\bs;\mu)
 =\sum_{N=0}^\infty h_N(\bs;\bdelta)
                  \CC_{\bq,A}(W+N;\mu),\qquad W=\sum_js_j.}
\end{equation}
The series is normally convergent on compact subsets where a strict
version of \eqref{clc:eq:rho-condition} holds, jointly in the independent
orders, the centers, and $\mu$ in a closed substrip of
$|\Ima\mu|<\pi\sum_jq_j^{-1}$. Every fixed mixed order or center
derivative may be taken term by term. Center collisions are regular.
If the scales are commensurate, every coefficient function $\CC$ on the
right has the finite pole-free Hurwitz or Lerch form already proved.
\end{theorem}
\begin{proof}
On the contour, $|p+A|\ge A+c$, so
$|\delta_j/(p+A)|\le\rho_0<1$ uniformly. The chosen logarithms satisfy
\[
 \log(p+A_j)=\log(p+A)+
             \log\left(1+\frac{\delta_j}{p+A}\right).
\]
This is first clear as $\delta_j$ moves from zero inside the strict disk,
and continues there because both $p+A$ and $p+A_j$ stay in the right
half-plane. The binomial series therefore gives
\[
 \prod_j(p+A_j)^{-s_j}
 =(p+A)^{-W}\sum_{N\ge0}h_N(\bs;\bdelta)(p+A)^{-N}.
\]
On a compact set of orders, $|(s)_n/n!|$ is at most a fixed polynomial
in $n$; this follows directly from the Gamma ratio away from a finite
initial set, with continuity covering nonpositive integer orders.
The absolute binomial series is consequently bounded by a convergent
polynomial-times-geometric series, uniformly along the whole contour.
The remaining factors in \eqref{clc:eq:U-contour} have an integrable
exponential vertical majorant, locally uniformly in all parameters.
Tonelli's theorem for the majorant and ordinary dominated convergence
justify termwise contour integration, giving \eqref{clc:eq:unequal-series}.

The same estimates hold on a slightly larger compact neighborhood of
any strict parameter set. Cauchy's integral formula then proves normal
convergence of every fixed derivative, not merely of the undifferentiated
series. In particular no singularity is created when two $\delta_j$
coincide. The final assertion follows by substituting
Theorem~\ref{clc:thm:finite} into each term.
\end{proof}

\begin{corollary}[Every common-strip configuration has such an expansion]
\label{clc:cor:global-center}
For any fixed centers satisfying \eqref{clc:eq:unequal-strip}, a reference
$A$ satisfying \eqref{clc:eq:rho-condition} exists. One may choose a real
number larger than zero and all the finitely many numbers
\begin{equation}\label{clc:eq:choose-A}
 \frac{|A_j|^2-c^2}{2(c+\Rea A_j)}.
\end{equation}
\end{corollary}
\begin{proof}
The denominators are positive. The strict inequality
$|A_j-A|^2<(A+c)^2$ is exactly
$|A_j|^2-c^2<2A(c+\Rea A_j)$, so the stated choice works for every $j$.
\end{proof}

This proves the unequal-shift research question in the preceding
reciprocal report \cite[Further research, question 2]{RLC} for its
right-half-plane shifts, and extends it to unequal scales whenever the
common strip exists. For equal scales $q_j=1$, the old shifts are
$a_j=1+A_j$; their conditions $\Rea a_j>0$ imply that such a common
$c<1$ can be chosen. Thus the original equal-scale question is fully
covered, including all independent spectral derivatives.

The result should not be described as a finite reduction of arbitrary
unequal centers: the outer series is generally infinite. The finite
closure is retained \emph{inside each term}. Theorem
\ref{clc:thm:alignment} below explains why generic unequal centers do
not recover pure order addition.

% END sections/10-unequal-centers.tex
% BEGIN sections/11-rigidity.tex
\section{Two precise rigidity statements}
\label{clc:sec:rigidity}

\subsection{When can the independent orders add purely?}
\begin{theorem}[Affine-center alignment is necessary and sufficient]
\label{clc:thm:alignment}
Fix positive scales and centers with a common strip. The convolution
function $\mu\mapsto\UU_{\bq,\bA}(\bs;\mu)$ depends on the orders only
through $W=\sum_js_j$, for all orders in a nonempty open set, if and
only if $A_1=\cdots=A_r$. For the \emph{unnormalized} profiles
$F_{1+A_j/q_j,s_j}(e^{q_jx})$, pure addition as a function of $\mu$
also requires $q_1=\cdots=q_r$.
\end{theorem}
\begin{proof}
Sufficiency in the normalized case is Theorem~\ref{clc:thm:addition}.
For necessity, weighted convolution is in $L^1$ as well as continuous,
so one can take its bilateral transform. If the function depends only
on $W$, its derivative in the direction
$\partial_{s_j}-\partial_{s_k}$ vanishes. By
\eqref{clc:eq:U-contour} and transform uniqueness, this gives
\[
 G_{\bq}(p)\prod_l(p+A_l)^{-s_l}
             [\log(p+A_k)-\log(p+A_j)]=0
\]
on the common vertical line. The prefactors never vanish there, so the
logarithms agree. Exponentiation implies $p+A_k=p+A_j$, hence
$A_k=A_j$. This holds for every pair.

In the unnormalized case there is an additional transform factor
$\prod_jq_j^{s_j}$. The same derivative now gives
$\log q_j-\log q_k-\log(p+A_j)+\log(p+A_k)=0$.
Thus $q_j(p+A_k)=q_k(p+A_j)$ on a line. Comparing the linear
coefficients in $p$ gives $q_j=q_k$, and then $A_j=A_k$.
\end{proof}

The statement concerns the entire convolution \emph{function} of $\mu$,
not a special numerical coincidence at one point or one order. It does
not rule out exceptional identities at $\mu=0$. Even with equal scales,
unequal centers already fail in the simplest example:
\begin{align*}
 \UU_{(1,1),(0,1)}((1,0);0)&=\zeta(2),\\
 \UU_{(1,1),(0,1)}((0,1);0)&=\zeta(2)-1.
\end{align*}
To verify these, the zero-order profile is independent of its center.
The two integrals therefore reduce respectively to the equal-center
weight-one formulas at $A=0$ and $A=1$.

\subsection{What commensurability really characterizes}
Consider the following explicitly restricted elementary class on $x<0$:
finite sums of functions
\begin{equation}\label{clc:eq:kernel-class}
 \frac{P_\nu(x)e^{\rho_\nu x}}
             {1-\epsilon_\nu e^{Qx}},
 \qquad Q>0,\quad |\epsilon_\nu|=1,
\end{equation}
where each $P_\nu$ is a polynomial, and the finite data $\rho_\nu$,
$\epsilon_\nu$ and $P_\nu$ are constant in $x$. The same $Q$ is used
for every term. Complex exponents are allowed.

\begin{theorem}[Common-step kernel rigidity]\label{clc:thm:commensurate}
For positive real scales $\bq$, the logistic convolution $K_{\bq}$
has a representation of the form \eqref{clc:eq:kernel-class} on $x<0$
if and only if all scale ratios are rational.
\end{theorem}
\begin{proof}
The sufficiency is \eqref{clc:eq:kernel}, since commensurate scales have
a positive common multiple $Q$.
For necessity, suppose such a finite representation exists and fix
$\delta>0$. On $x< -\delta$, expand the denominators as geometric
series. Integrating each term of its negative-half-line Laplace
transform gives meromorphic functions of $p$ whose poles lie in the
finite union
\[
 \bigcup_\nu\{\rho_\nu+nQ:n\ge0\}.
\]
Polynomial factors only increase pole multiplicities. The upper cutoff
$-\delta$ supplies exponential decay in $n$, so after omitting any
finite initial terms the integrated series is normally convergent on
compact subsets away from that pole set. This establishes the stated
meromorphic pole restriction, regardless of any cancellations between
terms of the finite representation.

The Laplace transform over $[-\delta,\infty)$ is holomorphic for
$\Rea p>0$. Indeed the weighted convolution estimates with arbitrarily
small positive weight show that $K_{\bq}(x)$ grows more slowly than
$e^{\varepsilon x}$ for every $\varepsilon>0$ as $x\to+\infty$.
The full transform agrees with $G_{\bq}(p)$ on $0<\Rea p<q_*$,
so meromorphic continuation forces every positive pole of $G_{\bq}$
to belong to the finite union above.

But $G_{\bq}$ has an uncancelled pole at every $nq_j$, $n\ge1$:
cosecant factors have no zeros, so another factor cannot remove it.
For a fixed $j$, two of the infinitely many $nq_j$ lie in the same
one of the finitely many progressions. Subtracting them gives
$(n-n')q_j=(k-k')Q$, hence $q_j/Q\in\mathbb Q$.
This holds for each $j$, proving commensurability.
\end{proof}

This theorem is a characterization of a fixed common-step elementary
kernel class. It is \emph{not} a claim that incommensurate integrals
cannot be expressed by other special functions, cannot have isolated
exact values, or yield arithmetically independent constants.
The next section supplies abundant isolated exact values for arbitrary
positive scales, emphasizing the distinction.

% END sections/11-rigidity.tex
% BEGIN sections/12-parity.tex
\section{Universal parity identities beyond commensurate scales}
\label{clc:sec:parity}

At negative integer total orders the resolvent power becomes a
polynomial. A local residue at zero then yields exact identities
without a common period. Put
\begin{equation}\label{clc:eq:Bq}
 B_{\bq}(z)=z^rG_{\bq}(z)
       =\prod_{j=1}^r\frac{\pi z/q_j}{\sin(\pi z/q_j)}.
\end{equation}
This is an even analytic function near zero with $B_{\bq}(0)=1$.

\begin{theorem}[Reflection-residue parity law]
\label{clc:thm:parity}
Let $q_j>0$ be arbitrary and let $N\ge0$ be an integer. Then
\begin{align}
 &\CC_{\bq,A}(-N;0)-(-1)^{r+N}\CC_{\bq,-A}(-N;0)\notag\\
 &\hspace{18mm}=[z^{r-1}](z+A)^N B_{\bq}(z).
 \label{clc:eq:parity}
\end{align}
The two ordinary integrals are covered simultaneously when
$|\Rea A|<q_*$. At these integer orders each side is a polynomial in
$A$, so the identity also gives their polynomial continuation to all $A$.
In particular, if $r+N$ is odd,
\begin{equation}\label{clc:eq:parity-central}
 \boxed{\displaystyle
 \CC_{\bq,0}(-N;0)=\frac12[z^{r-N-1}]B_{\bq}(z).}
\end{equation}
A coefficient with negative index is zero.
\end{theorem}
\begin{proof}
At $W=-N$, the contour integrand is
$G_{\bq}(p)(p+A)^N$. Choose $0<c<q_*$. At integer order the expression
is a polynomial in $A$, so the contour can be considered independently
of any branch of a resolvent power. Move the vertical contour from $c$
to $-c$. The only pole crossed is zero; horizontal edges vanish by
exponential decay. If $I_c(A)$ denotes the integral on $\Rea p=c$,
\[
 I_c(A)-I_{-c}(A)=\Res_{p=0}G_{\bq}(p)(p+A)^N.
\]
Since $G_{\bq}(-p)=(-1)^rG_{\bq}(p)$, changing variable $p\mapsto-p$
in the second vertical integral gives
$I_{-c}(A)=(-1)^{r+N}I_c(-A)$.
The residue on the right is exactly
$[z^{r-1}](z+A)^NB_{\bq}(z)$. On the common ordinary-integral domain,
$I_c(\pm A)=\CC_{\bq,\pm A}(-N;0)$, proving
\eqref{clc:eq:parity}; the polynomial continuation follows at once.
At $A=0$ and $r+N$ odd, the left side is twice the desired value,
giving \eqref{clc:eq:parity-central}.
\end{proof}

\subsection{Odd-depth logistic integrals}
For $r=2h+1$ and $N=0$, the theorem gives
\begin{equation}\label{clc:eq:odd-depth}
 K_{\bq}(0)=\frac12[z^{2h}]
                   \prod_{j=1}^{2h+1}\frac{\pi z/q_j}{\sin(\pi z/q_j)}.
\end{equation}
Every such value is a homogeneous polynomial in the inverse squared
scales with rational coefficients times $\pi^{2h}$.
To see both the rationality and the exact construction, use
\[
 \frac{u}{\sin u}=1+\frac{u^2}{6}+\frac{7u^4}{360}+\cdots,
\]
whose coefficients are rational by formal division of the sine series.
The first two nontrivial cases are
\begin{align}
 K_{(q_1,q_2,q_3)}(0)
   &=\frac{\pi^2}{12}\sum_{j=1}^3q_j^{-2},
 \label{clc:eq:universal3}\\
 K_{(q_1,\ldots,q_5)}(0)
   &=\frac{\pi^4}{720}
       \left(7\sum_{j=1}^5q_j^{-4}
               +10\sum_{i<j}q_i^{-2}q_j^{-2}\right).
 \label{clc:eq:universal5}
\end{align}
Equation \eqref{clc:eq:universal3} is the ordinary double integral
\eqref{clc:eq:intro-cubic}. For $(q_1,q_2,q_3)=(1,\sqrt2,\sqrt3)$,
for instance, its value is $11\pi^2/72$, despite the absence of a
common scale period. For $(1,2,3)$ it recovers $49\pi^2/432$ from the
explicit Laurent table.

More generally,
\begin{equation}\label{clc:eq:negative-examples}
 \CC_{(q_1,q_2),0}(-1;0)=\frac12,\qquad
 \CC_{(q_1,q_2,q_3),0}(-2;0)=\frac12.
\end{equation}
If $r+N$ is odd and $N>r-1$, the value is exactly zero.
Such identities involve negative-order Lerch or polylogarithmic
profiles, which can change sign, and so do not contradict positivity
of the zero-order logistic convolution.

At nonzero $A$, equation \eqref{clc:eq:parity} supplies a finite
reflection identity instead of an evaluation of each side separately.
For example, at $r=2$, $N=2$ the residue is $2A$, so
\[
 \CC_{(q_1,q_2),A}(-2;0)-\CC_{(q_1,q_2),-A}(-2;0)=2A.
\]
This exact reflection is independent of the two scales. The theorem
makes no assertion that the complementary-parity central values have
a comparable universal polynomial form.

% END sections/12-parity.tex
% BEGIN sections/13-audit.tex
\section{Proof status, correction safeguards, and reproducible checks}
\label{clc:sec:audit}

\subsection{What was actually inspected}
The canonical manuscript README and its source-directory structure were
retrieved through the connected GitHub repository. The branch commit
recorded during this audit is
\begin{center}\small
\src{338a3a3877c337f981072061c16bc8ee1d87155d}.
\end{center}
The repository's timestamp for that commit is 11 October 2026 at
01:59:45 UTC, which is 10 October in the user's Pacific timezone.
The full accessible TeX source of \emph{Reciprocal Lerch Convolutions}
was retrieved from the user's Library and preserved locally for this
research pass. Its profile definitions, transform and weighted-holomorphy
proof, equal-scale closure, moment formulas, Stieltjes conventions,
harmonic-remainder identity, and future questions are the relevant
source material. That report records its own, earlier repository pin;
it must not be confused with the pin above.

The current incoming directory listing includes the bilinear Mellin,
cubic mixed-order, exact-resonance, independent-order, harmonic parity,
harmonic shift, higher reflection, nonlinear transport, ordered Hurwitz,
and unequal-twist packages. Their filenames and GitHub metadata were
visible, but their binary ZIP members were not read in this pass.
Accordingly this article does not assert that its results are absent
from every incoming archive, and does not claim to have audited those
archives. The source-audit artifact records this limitation explicitly.
A filename alone is not mathematical evidence of either novelty or overlap.

\subsection{Corrections and integration safeguards}
No new error in the inspected equal-scale theorem is asserted here.
The following are concrete corrections to plausible \emph{extensions}
and implementation mistakes, supported by exact formulas in this article.

\paragraph{Scale normalization.}
Replace a raw order-addition claim for unequal scales by the normalized
profiles \eqref{clc:eq:h}. The exact pair \eqref{clc:eq:swap} refutes the
unnormalized claim. Alignment means $q_j(a_j-1)=A$, not equality of the
original Lerch shifts $a_j$.

\paragraph{Pole cancellation before specialization.}
Store $\EE_a(v)$ as an entire object with value one at zero, and store
simple-pole Hurwitz sums as differences. The $1/3$ in
\eqref{clc:eq:zero-examples} is lost by zero-times-pole substitution.
In particular the expression $W\zeta(W+1)$ is not zero at $W=0$.

\paragraph{Periods and branches.}
The pole period is a common multiple of the positive scales.
Using a common denominator instead generally misidentifies the pole
lattice. At positive $\mu$ in the periodic case, use the entire combined
Lerch expression; its zero jump was proved, not assumed. Independently
choosing branch values and discarding imaginary parts is not the rule.

\paragraph{Independent orders and common strips.}
An integer-only partial fraction formula cannot be differentiated in
its changing integer summation limit. The normal expansion in
Theorem~\ref{clc:thm:unequal} supplies an actual holomorphic family.
For unequal scales and centers, the common-strip hypothesis
\eqref{clc:eq:unequal-strip} is explicit; separate one-profile
admissibility is not silently substituted for it.

\paragraph{Primitive constants and residue scope.}
Use anchored Taylor remainders for higher primitives. Their resonant
quotients have specified removable values. The universal parity result
is an equality of the stated ordinary integrals on their common domain,
then a polynomial continuation at negative integer order; it is not a
license to continue arbitrary complex-order integrals without tracking
branches and convergence.

\paragraph{Arithmetic and numerical scope.}
A functional pole-lattice obstruction is not arithmetic independence.
The current $S_6$ and revised $S_8$ statuses are unchanged. During
exploratory numerical work, an ordinary Python expression $3^{-2}$
introduced a binary floating-point coefficient into an otherwise
high-precision comparison; replacing it by exact or multiprecision
arithmetic removed that discrepancy. The delivered numerical routines
convert rational constants without such a binary-float intermediate.
This was an implementation issue in the present exploratory work, not
an error attributed to the repository.

\subsection{Executed verification}
The delivered completed runs comprise \textbf{307 exact symbolic
assertions, 91 numerical diagnostics, and five negative controls}.
The exact and numerical roles are deliberately separate.
\begin{center}
\begin{tabular}{l r r r}
\toprule
Script/run & Exact & Numerical & Negative controls\\
\midrule
\texttt{verify.py --part exact} & 285 & 0 & 2\\
\texttt{verify.py --part numeric} & 0 & 77 & 2\\
\texttt{verify\_extra.py} & 22 & 14 & 1\\
\midrule
Total & 307 & 91 & 5\\
\bottomrule
\end{tabular}
\end{center}
The exact suite compares the Bell and direct reciprocal-sine algorithms
for eighteen scale families, checks period residues and scale covariance,
verifies the full printed $(1,2,3)$ table, and checks the two-scale pole
formulas, mixed-moment recurrence, harmonic Pochhammer coefficients,
primitive resonance denominators, and low-depth universal parity
polynomials. The corruption control changes an actual coefficient of
the Laurent table and verifies that the independent construction rejects
it. A separate domain control rejects a zero scale.

The numerical suite compares finite central closure with independent
vertical-contour quadrature, including complex orders, complex centers,
zero order, positive order, and negative-integer resonance. It also
compares noncentral Lerch closure, ordinary logarithmic product integrals,
weighted moments, Gamma primitives, Stieltjes jets, balanced harmonic
remainders, unequal-center expansions, and irrational-scale parity
values. Higher primitive tests compare the closed formulas to
32-point Gauss--Legendre integration; their resonant order derivatives
are evaluated by a separate 48-point Cauchy sum. The universal cubic
formula is also checked against ordinary one-dimensional integrals
after an elementary two-factor convolution.

The usual working precision is 55 decimal digits. The balanced-remainder
integral uses 100 digits to control cancellation and a truncated
real-line integration interval; its acceptance tolerance is $10^{-30}$.
The other default tolerances are $10^{-36}$, with $10^{-42}$ for the
unequal-center expansion tests. The finite quadrature and finite expansion
errors are not interval-enclosed. An initial adaptive higher-primitive
diagnostic exceeded the execution limit; the delivered fixed quadrature
and Cauchy-sum checks completed successfully. All completed records,
including actual residuals and software versions, are in
\texttt{verification/*-results.json}.

These computations test the formulas and their implementation. They are
not interval certificates, not proofs for all parameters, and not a
proof-assistant formalization. The all-index and continuum assertions
rest on the analytic proofs above.

\subsection{Integration status}
\begin{center}
\begin{tabular}{p{0.56\textwidth}p{0.35\textwidth}}
\toprule
Claim & Status in this delivery\\
\midrule
Commensurate-scale, all-depth, entire-order finite closure & Proved.\\
Every mixed logarithmic moment and independent order jet & Proved by a finite algorithm.\\
Unequal-center expansion in the common-strip domain & Proved, normally convergent.\\
All anchored primitive orders and their resonances & Proved.\\
Universal negative-integer parity for arbitrary positive scales & Proved.\\
Pure-order-addition and common-step kernel rigidity & Proved in the stated functional classes.\\
Canonical Gaussian $S_6$ and revised $S_8$ & Not proved or disproved here.\\
Novelty relative to every unread incoming ZIP & Not established.\\
Historical priority for every special case & Not claimed.\\
\bottomrule
\end{tabular}
\end{center}
All theorem and equation labels use the prefix \texttt{clc:} to reduce
integration collisions. The package includes modular sources, a standalone
source, a compiled PDF, reusable verification code, provenance notes,
and a checksum manifest. No change to the canonical repository was made.

% END sections/13-audit.tex
% BEGIN sections/14-questions.tex
\section{Further research questions and concrete next targets}
\label{clc:sec:questions}

The following questions concern further exact identities rather than
inequalities or asymptotic error bounds. They also distinguish genuine
remaining problems from the two continuation questions resolved here.

\begin{enumerate}[label=\textbf{Q\arabic*.},leftmargin=2.8em,itemsep=9pt]
\item \textbf{A compact multivariate tilt generating function.}
For the constraint $\sum_jx_j=\mu$, introduce independent factors
$e^{\lambda_jx_j}$. Their transform becomes
$\prod_jg_{q_j}(p-\lambda_j)(p+A-\lambda_j)^{-s_j}$.
Find a compact representation in a specified multivariate Lerch or
Lauricella family, with every collision and normalization explicit,
that avoids an outer center-expansion series. The present moment theorem
already determines every Taylor coefficient, but does not supply such
a compact resummation. A characterization of its genuinely finite
sectors would be more informative than an unqualified ``closed form.''

\item \textbf{The complementary negative-order parity sector.}
Theorem~\ref{clc:thm:parity} evaluates the sector $r+N$ odd at $A=0$.
For $r+N$ even and incommensurate scales, seek exact representations in
Barnes-type or multiple Hurwitz functions that retain the full scale
symmetry. A correct formula should reduce to the finite commensurate
closure when all scale ratios become rational, without assuming that
the limiting pole multiplicities are unchanged.

\item \textbf{Minimal cyclotomic coordinates and conductor descent.}
The finite tables live in explicit real cyclotomic algebras, but those
algebras and the Hurwitz argument grids are often redundant. Determine
an exact reduction algorithm for the coefficient field and the
Hurwitz-coordinate conductor simultaneously. The target is a canonical
functional reduction with exact certificates, not numerical independence
of a chosen set of periods. The two-scale beta and character-three
formulas are required base cases.

\item \textbf{Finite closure without pure order addition.}
Alignment is characterized exactly for pure addition. It does not
classify all finite identities with unequal centers. Positive integer
orders permit finite partial fractions, while our arbitrary-order
series is usually infinite. Classify exceptional affine-center and
order subvarieties on which the outer series terminates or resums
finitely. Nonpositive integer orders, repeated centers, and rational
center differences provide distinct test cases and should not be
conflated.

\item \textbf{Beyond a single common-step elementary kernel class.}
Theorem~\ref{clc:thm:commensurate} allows one shared positive $Q$.
Determine the exact obstruction when finitely many different denominator
periods are allowed in the elementary representation. The uncancelled
pole support alone may no longer suffice; the Laurent coefficient
sequences along intersecting progressions carry additional information.
Any stronger theorem should state its representational class precisely,
and should remain compatible with the universal isolated parity values.

\item \textbf{Several independent multiplicative constraints.}
Replace the single linear condition on logarithms by an integer matrix
of constraints. The transform is then a product of cosecants and affine
resolvents on a higher-dimensional space. Seek a finite exact
polylogarithm/Hurwitz closure when the pole hyperplane arrangement is
rational and the resolvents are aligned along the appropriate dual
subspaces. A first nontrivial target is two independent constraints
with four or five factors, including intersecting-pole residues and all
independent order derivatives.

\item \textbf{Colored profiles and exact contour-crossing terms.}
Replace the negative argument of the Lerch profile by a root-of-unity
multiple, and allow complex logarithmic shifts. Derive the correct
weighted strip and the finite residue or branch terms when a pole
crosses the chosen contour. The jump cancellation of
Theorem~\ref{clc:thm:finite} is a special aligned case, not a general
rule that all colored jumps vanish. Exact quarter-turn specializations
could connect this extension to the manuscript's Gaussian identities.

\item \textbf{A canonical finite basis for mixed moments.}
The algorithm in Theorem~\ref{clc:thm:moments} is terminating but not
claimed minimal. Determine all functional syzygies among its differentiated
Laurent tables and produce a canonical reduced representative that
respects center differentiation, scale covariance, and
\eqref{clc:eq:moment-sum}. The equal-scale central-factorial recurrence
is an instructive subcase; unequal periods introduce new redundancies.

\item \textbf{Machine-checked analytic infrastructure.}
Formalize weighted entire families, bilateral convolution inversion,
and the finite period-residue argument in a proof assistant. The exact
coefficient tables and negative controls can be used as finite fixtures,
but they do not replace a formal proof of contour deformation, normal
convergence, or branch cancellation. The first useful formal milestone
is the all-complex-order two-scale theorem with its removable resonances.
\end{enumerate}

The commensurate-dilation problem and the unequal-shift order-jet
expansion posed in \cite{RLC} now have complete proofs in the stated
settings. The next advances should go beyond those solved settings,
while preserving their normalization and convergence checks. In
particular, connecting these integral families to the unresolved
canonical Gaussian reduction vectors would require a separate explicit
identity or exact reduction certificate; it is not a consequence of
finite closure alone.

% END sections/14-questions.tex
% BEGIN references.tex
\begingroup
\raggedright
\begin{thebibliography}{9}
\addcontentsline{toc}{section}{References}
\bibitem{Repo}
ProveIt Contributors,
\emph{Polylogarithms and their Arithmetic Bridges}, canonical manuscript
README and source structure, VladimirReshetnikov/ProveIt.
Repository commit recorded during this audit:
\src{338a3a3877c337f981072061c16bc8ee1d87155d}.
\url{https://github.com/VladimirReshetnikov/ProveIt/tree/main/Analysis/Polylogarithms/docs/manuscript}.
Incoming metadata:
\url{https://github.com/VladimirReshetnikov/ProveIt/tree/main/docs/incoming}.
Accessed 10 October 2026, Pacific time.

\bibitem{RLC}
ProveIt research continuation,
\emph{Reciprocal Lerch Convolutions: Entire-order identities,
central-factorial zeta reductions, all logarithmic moments, and
Stieltjes--Gamma calculus}, 10 October 2026.
User Library source:
\src{ProveIt_Reciprocal_Lerch_Convolutions_2026-10-10.tex}.
The report records the earlier repository anchor
\src{0c740647c48baec7f0b4192723e977d69917dd52}.
In particular, see its weighted transform, equal-scale closure, and
further research questions 2 and 5. The source hash is retained in the
accompanying source-audit artifact.

\bibitem{DLMF25}
F. W. J. Olver et al., eds.,
\emph{NIST Digital Library of Mathematical Functions}, Chapter 25,
Sections 25.11 (Hurwitz zeta), 25.12 (polylogarithms), and 25.14
(Lerch transcendent).
\url{https://dlmf.nist.gov/25.11},
\url{https://dlmf.nist.gov/25.12},
\url{https://dlmf.nist.gov/25.14}.
The standard definitions, integral representations, and
$\zeta'(0,a)$ normalization were consulted for this work.

\bibitem{DLMF5}
F. W. J. Olver et al., eds.,
\emph{NIST Digital Library of Mathematical Functions}, Chapter 5,
Sections 5.5 (Gamma functional relations) and 5.15 (polygamma functions).
\url{https://dlmf.nist.gov/5.5},
\url{https://dlmf.nist.gov/5.15}.

\bibitem{SS}
J. P. Selvaggi and J. A. Selvaggi,
``The Application of Real Convolution for Analytically Evaluating
Fermi-Dirac-Type and Bose-Einstein-Type Integrals,''
\emph{Journal of Complex Analysis} \textbf{2018} (2018), article 5941485.
\url{https://doi.org/10.1155/2018/5941485}.
Cited for the prior use of real convolution in this class of integrals,
not as a source of the present all-scale theorems.
\end{thebibliography}
\endgroup

% END references.tex
\end{document}
